\documentclass[11pt,a4paper]{article}
\usepackage[T1]{fontenc}
\usepackage[utf8]{inputenc}
\usepackage[spanish,es-noshorthands]{babel}
\usepackage{amsmath,amssymb}
\usepackage{graphicx,booktabs,hyperref,listings,geometry,enumitem}
\geometry{margin=2.4cm}
\hypersetup{colorlinks=true,linkcolor=black,urlcolor=blue}
\lstset{
  basicstyle=\ttfamily\footnotesize,
  breaklines=true,
  columns=fullflexible,
  frame=single,
  framesep=3pt,
  showstringspaces=false,
  keepspaces=true,
  language=bash,
}
\setlist[itemize]{topsep=0.3em,itemsep=0.15em}
\setlist[enumerate]{topsep=0.3em,itemsep=0.15em}
\setlength{\emergencystretch}{12em}

\title{Arquitectura, flujo y ejecución del proyecto\\
\texttt{fibonacci\_photonics}}
\author{Emanuel Orozco Gallego}
\date{28 de septiembre de 2026}

\begin{document}
\maketitle

\begin{abstract}
Manual de arquitectura del repositorio que reproduce
Reyes-Gómez \textit{et al.}, Phys.\ Rev.\ B \textbf{81}, 153101 (2010).
Describe cómo está organizado el código, cómo circulan los datos desde un
YAML hasta una figura, qué dependencias hacen falta y cómo instalar y
ejecutar el proyecto de punta a punta.
La física está en \textit{Guía conceptual}; el algoritmo TMM, en
\textit{Solución numérica e implementación}; el método de ondas planas, en
\textit{El método de expansión en ondas planas}.
\end{abstract}

\tableofcontents
\newpage

\section{Qué es este repositorio}

Es una reimplementación en Python del Brief Report de 2010 con dos métodos
independientes: la matriz de transferencia (TMM), que es el método de
producción, y la expansión en ondas planas (PWE), que la verifica. No hay
servidor web ni interfaz gráfica: el producto son \emph{datos}, \emph{figuras}
y \emph{PDFs}, separados por método en \texttt{results/tmm/},
\texttt{results/pwe/} y \texttt{results/comparison/}.

Tres capas, de dentro hacia afuera:

\begin{enumerate}
\item \textbf{Biblioteca} \texttt{src/fibonacci\_photonics/}: física y numérica.
No sabe de Matplotlib ni de números de figura.
\item \textbf{Scripts} \texttt{scripts/\{tmm,pwe,comparison\}/}: leen YAML, llaman a la biblioteca,
dibujan y escriben archivos.
\item \textbf{Documentos} \texttt{docs/}: LaTeX que incrusta las figuras
generadas.
\end{enumerate}

Los parámetros científicos no viven en el código: viven en
\texttt{configs/figure\_0N.yaml}. Cambiar un plasma o un ángulo no exige
editar Python.

\section{Árbol del proyecto}

\begin{lstlisting}
analisis_numerico/
  src/fibonacci_photonics/   biblioteca (layout src)
    core/                    constantes, materiales, Fibonacci, parametros, bandas
    tmm/                     matriz de transferencia y barrido
    pwe/                     ondas planas: Fourier, solver k(w), forma w(k)
    analysis/                modos, anchos, comparacion, bordes de banda
    plotting.py              estilo y guardado de figuras
  src/fibonacci_tmm/         alias del nombre antiguo (compatibilidad)
  tests/{core,tmm,pwe,analysis,compat}/   pytest
  configs/                   YAML por figura (fisica) + base.yaml
  configs/pwe/               parametros numericos del PWE
  scripts/_common.py         utilidades compartidas
  scripts/tmm/               Figs. 1-6 y convergencia con TMM
  scripts/pwe/               benchmark del libro, estudios y Figs. 1-6 con PWE
  scripts/comparison/        superposiciones, errores, tablas y tiempos
  scripts/reproduce_*.py     envoltorios antiguos (llaman a scripts/tmm/)
  scripts/verify_results.py  verificacion automatica (make verify)
  results/tmm/figure_0N/     data/, output/, README.md
  results/pwe/<estudio>/     data/, output/
  results/comparison/        figuras, summary.json y tables/*.tex
  figures -> results/tmm     enlace de compatibilidad
  references/pwe/            PWE-1d.ipynb, PWE_2D.m (el PDF del libro es local)
  references/paper/          articulo APS y recortes (local, no se sube)
  docs/                      auditoria, papers, guias, este documento
  presentacion/              presentacion (no la toca make results)
  notebooks/                 exploracion
  pyproject.toml             paquete, pytest, pythonpath
  requirements.txt           pip
  Makefile                   tests, resultados, PDFs
\end{lstlisting}

El entorno virtual \texttt{.venv/} y las cachés (\texttt{\_\_pycache\_\_},
\texttt{.pytest\_cache}) no forman parte del código fuente.
\texttt{.gitignore} también excluye auxiliares de LaTeX y el directorio
\texttt{\_paper\_pages/} (recortes del PDF original).

\section{Arquitectura de la biblioteca}

El paquete usa el \emph{src layout}: las importaciones son
\texttt{from fibonacci\_photonics.\ldots} y Setuptools busca paquetes en
\texttt{src/} (\texttt{pyproject.toml}).
Pytest añade \texttt{src} a \texttt{PYTHONPATH}. Los scripts hacen
\texttt{sys.path.insert} a \texttt{src/} al arrancar.

El nombre antiguo \texttt{fibonacci\_tmm} se conserva como alias: cada módulo
\texttt{fibonacci\_tmm.X} reemplaza su entrada en \texttt{sys.modules} por el
módulo nuevo correspondiente, así que los cuadernos, la presentación y
cualquier código externo siguen funcionando sin cambios
(\texttt{tests/compat/test\_legacy\_imports.py} lo comprueba).

\subsection{Módulos y responsabilidades}

{\footnotesize
\begin{center}
\begin{tabular}{llp{6.4cm}}
\toprule
Módulo & Depende de & Responsabilidad \\
\midrule
\texttt{core/constants.py} & --- & \(c\), GHz, mm, tolerancias \\
\texttt{core/fibonacci.py} & constants & \(S_m\), \(F_m\), \(L_m\) \\
\texttt{core/materials.py} & --- & aire y Drude \(\varepsilon(\omega),\mu(\omega)\) \\
\texttt{core/electromagnetics.py} & constants & \(n\), \(q\), \(Q\), \(\chi\), TE/TM \\
\texttt{core/params.py} & materials, electromagnetics & YAML \(\to\) SI \\
\texttt{core/bands.py} & --- & \texttt{DispersionScan}, máscara de banda, \(kL_m/\pi\) \\
\texttt{core/effective\_medium.py} & materials & \(\langle n\rangle\), gap de índice nulo \\
\texttt{core/provenance.py} & params, constants & metadatos JSON de cada corrida \\
\texttt{tmm/transfer\_matrix.py} & core & \(M\), \(T_m\), \(R_m\) \\
\texttt{tmm/dispersion.py} & transfer\_matrix, bands & barrido TMM \\
\texttt{pwe/fourier.py} & fibonacci & coeficientes de la celda, Toeplitz \\
\texttt{pwe/bloch.py} & fourier, params & QEP en \(k\), selección del modo \\
\texttt{pwe/dispersion.py} & bloch, bands & barrido PWE, \texttt{SemitraceEvaluator} \\
\texttt{pwe/eigenfrequency.py} & fourier & forma \(\omega(k)\) del libro y de Drude \\
\texttt{analysis/plasmon\_modes.py} & bands, fibonacci & recuento de subbandas \\
\texttt{analysis/bandwidth.py} & plasmon\_modes, band\_edges & \(\nu_{\min},\nu_{\max}\) vs \(\theta\) \\
\texttt{analysis/comparison.py} & bands & métricas TMM vs PWE, Brent \\
\texttt{analysis/band\_edges.py} & comparison & localización, refinamiento y validación de subbandas \\
\texttt{plotting.py} & bands, matplotlib & estilo PRB, curvas, \texttt{save\_figure} \\
\texttt{\_\_init\_\_.py} & los anteriores & API pública \\
\bottomrule
\end{tabular}
\end{center}
}

La dependencia es acíclica y va de lo más físico a lo más aplicado:
\texttt{core} \(\to\) \texttt{tmm} y \texttt{pwe} (en paralelo, sin depender
uno del otro) \(\to\) \texttt{analysis} \(\to\) dibujo. Los dos métodos
producen el mismo objeto \texttt{DispersionScan}, de modo que el análisis
(modos, anchos, bordes) y el dibujo no saben qué método generó los datos.

\texttt{plotting.py} es el único módulo de la biblioteca que importa
Matplotlib. El núcleo numérico puede usarse sin pintar nada.

\subsection{Objetos que cruzan las fronteras}

\begin{itemize}
\item \texttt{SuperlatticeSpec}: espesores en metros y Drude en rad/s.
Lo construye \texttt{spec\_from\_mapping} a partir del YAML.
\item \texttt{DispersionScan}: vectores \(\nu\), \(R_m\), máscara de banda,
\(k L_m/\pi\). Lo produce \texttt{scan\_dispersion}.
\item \texttt{FrequencyInterval}: un trozo conexo \([\nu_{\min},\nu_{\max}]\).
\item \texttt{BandwidthPoint}: anchos y centros a un \(\theta\) dado (Fig.~6).
\end{itemize}

\section{Capa de configuración}

\begin{center}
\begin{tabular}{ll}
\toprule
Archivo & Uso \\
\midrule
\texttt{configs/base.yaml} & \(a=b=12\,\mathrm{mm}\), aire, \(c\) SI \\
\texttt{configs/figure\_01.yaml} & \(\nu_e=\nu_m=3\,\mathrm{GHz}\), \(m=3,4\), cuatro \(\theta\) \\
\texttt{configs/figure\_02.yaml} & zoom \(2\)--\(4\,\mathrm{GHz}\), \(m=3\ldots 6\) \\
\texttt{configs/figure\_03.yaml} & como Fig.~1 con \(\nu_m=1\,\mathrm{GHz}\) \\
\texttt{configs/figure\_04.yaml} & ventanas densas cerca de \(1\,\mathrm{GHz}\) \\
\texttt{configs/figure\_05.yaml} & palos vs \(m=2\ldots 8\) \\
\texttt{configs/figure\_06.yaml} & \(\nu(\theta)\) relleno, \(m=2\ldots 7\) (inferida) \\
\texttt{configs/pwe/figure\_0N.yaml} & armónicos por capa, mallas y órdenes del PWE \\
\texttt{configs/pwe/book\_benchmark.yaml} & cristal \(\varepsilon_1=1\), \(\varepsilon_2=9\) del libro \\
\texttt{configs/pwe/convergence.yaml} & estudios de convergencia y costo \\
\texttt{configs/pwe/spectral\_pollution.yaml} & forma \(\omega(k)\) con Drude \\
\texttt{configs/comparison/efficiency.yaml} & estudio de tiempos TMM vs PWE en un núcleo \\
\bottomrule
\end{tabular}
\end{center}

La física de cada figura está en un solo sitio (\texttt{configs/figure\_0N.yaml})
y la leen los dos métodos; \texttt{configs/pwe/} solo añade parámetros
numéricos del PWE.

Cada YAML mezcla \emph{física} (espesores, plasmas, polarización) con
\emph{numérica de gráfica} (malla de \(\nu\), lista de \(\theta\), ventanas).
\texttt{load\_spec} / \texttt{spec\_from\_mapping} solo leen la física.
El resto lo consume el script de la figura.

Campo \texttt{source}: \texttt{paper} si el original da los números;
\texttt{inferred\_from\_figures\_3\_to\_5} en la Fig.~6.

\section{Capa de scripts}

\texttt{scripts/\_common.py} reúne utilidades compartidas: cargar YAML,
crear los directorios \texttt{results/<método>/<estudio>/\{data,output\}},
guardar archivos \texttt{.npz} y JSON, escribir el sidecar de provenance y
repartir tareas entre procesos (\texttt{parallel\_map}, con un hilo de BLAS
por proceso; \texttt{PWE\_WORKERS=n} limita su número).

Cada \texttt{scripts/tmm/figure\_0N.py} hace el mismo esqueleto:

\begin{enumerate}
\item \texttt{load\_figure\_config("figure\_0N.yaml")}.
\item Construir la malla de \(\nu\) (y de \(\theta\) si aplica).
\item Llamar a \texttt{scan\_dispersion} o a \texttt{bandwidth\_versus\_angle}.
\item Dibujar con \texttt{plotting}.
\item Escribir PDF, PNG (300 dpi), SVG, datos y \texttt{README.md} en
\texttt{results/tmm/figure\_0N/}.
\end{enumerate}

\texttt{scripts/pwe/figure\_0N.py} sigue el mismo esqueleto con
\texttt{fibonacci\_photonics.pwe.scan\_dispersion} y escribe en
\texttt{results/pwe/figure\_0N/}. Además, \texttt{scripts/pwe/} contiene el
benchmark del libro (\texttt{book\_benchmark.py}), la ilustración de la celda
(\texttt{cell\_fourier.py}), la contaminación espectral de la forma
\(\omega(k)\) (\texttt{spectral\_pollution.py}) y la convergencia y el costo
(\texttt{convergence.py}).

\texttt{scripts/comparison/} lee los datos de los dos métodos, recalcula la TMM
en la misma malla que el PWE y escribe superposiciones, errores,
\texttt{summary.json} y las tablas LaTeX de \texttt{results/comparison/tables/}
que incluyen el paper del PWE y el anexo de comparación.
\texttt{efficiency.py} mide en un núcleo el tiempo por frecuencia frente a \(m\),
el diagrama trabajo-precisión, el costo de cada figura y la memoria
(\texttt{--plot-only} redibuja sin medir).

Cada carpeta tiene un \texttt{run\_all.py} que lanza sus scripts en orden con
el mismo intérprete (\texttt{sys.executable}) y aborta si uno falla
(\texttt{check=True}). Los envoltorios \texttt{scripts/reproduce\_figure\_0N.py}
y \texttt{scripts/reproduce\_all.py} ejecutan los de \texttt{scripts/tmm/} con
\texttt{runpy}.

\texttt{scripts/tmm/convergence.py} no pinta el paper: estudia cómo
\(\Delta\nu\) de la Fig.~4 cambia con el número de puntos.

\section{Flujo de una corrida}

El camino feliz, para \emph{una} figura, es:

\begin{lstlisting}
configs/figure_0N.yaml
        |
        v
  spec_from_mapping          (mm, GHz -> SI)
        |
        v
  SuperlatticeSpec
        |
        +--> tmm.scan_dispersion(spec, m, theta, nu, TE)
        |         +--> semitrace_by_recurrence   (por defecto)
        +--> pwe.scan_dispersion(spec, m, theta, nu, TE, n_max)
        |         +--> QEP en k (regla inversa), R = Re cos(kL)
        |         v
        |    DispersionScan   (mismo objeto para los dos metodos)
        |
        +--> detect_plasmon_modes / allowed_intervals   (Figs. 2, 4, 5)
        +--> bandwidth_versus_angle                     (Fig. 6)
        |
        v
  plotting.save_figure
        |
        +--> results/<metodo>/figure_0N/output/figure_0N.{pdf,png,svg}
        +--> results/<metodo>/figure_0N/data/*.npz
        +--> results/<metodo>/figure_0N/data/summary.json
        +--> results/<metodo>/figure_0N/data/metadata.json
\end{lstlisting}

Los PDFs científicos no recalculan física: incluyen las figuras ya generadas
en \texttt{results/} (o en \texttt{figures/}, que apunta a
\texttt{results/tmm/}). Por eso hay que generar los resultados \emph{antes} de
compilar los documentos si se cambió el código.

Flujo del \texttt{Makefile}:

\begin{lstlisting}
make all
  make tests      -> pytest tests
  make results    -> make tmm pwe comparison
    make tmm        -> scripts/tmm/run_all.py (Figs. 1-6 + convergencia)
    make pwe        -> scripts/pwe/run_all.py (horas de CPU)
    make comparison -> scripts/comparison/run_all.py (incluye eficiencia)
  make documents  -> pwe-docs, guides, paper, docs
    make pwe-docs   -> metodo_ondas_planas, comparacion_tmm_pwe
    make guides     -> guia_conceptual, solucion_numerica, arquitectura
    make paper      -> paper_reimplementation
    make docs       -> software_documentation
  make verify     -> scripts/verify_results.py
make reproduce    -> results + documents + verify (sin tests)
make figures      -> alias de make tmm
make presentation -> presentacion/ (objetivo aparte, no lo llama make all)
\end{lstlisting}

Todos los PDF se compilan con una sola regla de patrón (dos pases de
\texttt{pdflatex}, que se detiene en el primer error). Cada PDF depende de su
\texttt{.tex} y de las figuras y tablas que incluye, así que \texttt{make}
solo recompila lo que cambió. El PDF se copia también a la raíz del repo
(esas copias no se suben) para abrirlo sin entrar en \texttt{docs/}.

\texttt{make verify} comprueba que existan todas las salidas, que se cumplan los
criterios de acuerdo TMM vs PWE (conteos $F_{m-2}$, errores de semitraza, de
borde y de ancho), que cada tabla y figura citada exista y que cada PDF esté al
día y compile sin errores ni referencias indefinidas. Con
\texttt{--reference DIR} compara además los datos con una corrida anterior.

\section{Tests}

\texttt{pyproject.toml} configura \texttt{testpaths = ["tests"]} y
\texttt{pythonpath = ["src"]}. No hace falta instalar el paquete para
testear.

\begin{center}
\begin{tabular}{ll}
\toprule
Archivo & Qué garantiza \\
\midrule
\texttt{core/test\_fibonacci.py} & palabras, inflación, \(N_B\), \(L_m\) \\
\texttt{core/test\_materials.py} & ceros de Drude, \(\nu_e\), \(\nu_m\) \\
\texttt{core/test\_params.py} & YAML de las seis figuras \\
\texttt{tmm/test\_transfer\_matrix.py} & \(\det T=1\), eqs.\ (7)--(9), dos vías de \(R_m\) \\
\texttt{tmm/test\_dispersion.py} & máscara de banda, gap \(\langle n\rangle=0\) \\
\texttt{pwe/test\_fourier.py} & ec.~(4.38), fracción de llenado, Toeplitz, síntesis \\
\texttt{pwe/test\_solver.py} & libro, \(R_{\rm PWE}\approx R_{\rm TMM}\), Laurent vs inversa, contaminación \\
\texttt{analysis/test\_validation.py} & \(N_{\mathrm{modos}}=F_{m-2}\) \\
\texttt{analysis/test\_comparison.py} & métricas, Brent, localización y validación de bandas \\
\texttt{compat/test\_legacy\_imports.py} & el alias \texttt{fibonacci\_tmm} \\
\bottomrule
\end{tabular}
\end{center}

\texttt{conftest.py} fija un hilo de BLAS (las matrices del PWE son pequeñas y
el multihilo solo añade sobrecarga) y expone dos fixtures de \texttt{SuperlatticeSpec}:
plasmas iguales (\(3,3\,\mathrm{GHz}\), Figs.~1--2) y partidos
(\(3,1\,\mathrm{GHz}\), Figs.~3--6).

Tres casos se omiten (\texttt{skip}) cuando una identidad de Fibonacci pide
\(k>m\). El resto debe pasar.

\section{Dependencias}

\subsection{Tiempo de ejecución (Python)}

Declaradas en \texttt{requirements.txt} y en \texttt{pyproject.toml}.
Python \(\ge 3.10\) (\texttt{requires-python}). El desarrollo se hizo con
Python 3.14.

\begin{center}
\begin{tabular}{lll}
\toprule
Paquete & Mínimo & Para qué \\
\midrule
\texttt{numpy} & 1.24 & arrays, \(\sqrt{\,}\), trigonométricas complejas \\
\texttt{scipy} & 1.10 & Brent, \texttt{eigh} generalizado del PWE \\
\texttt{matplotlib} & 3.7 & Figs.~1--6 \\
\texttt{pyyaml} & 6.0 & configs \\
\texttt{pytest} & 7.0 & tests (extra \texttt{dev} en pyproject) \\
\bottomrule
\end{tabular}
\end{center}

No hay CUDA, MPI ni bases de datos. No se descarga nada en tiempo de
cálculo: el paper original es un PDF local.

\subsection{Sistema, para los PDFs}

Además de Python:

\begin{itemize}
\item \texttt{make}
\item \texttt{pdflatex} (TeX Live). Estilos usados: \texttt{article},
\texttt{babel-spanish}, \texttt{amsmath}, \texttt{graphicx},
\texttt{booktabs}, \texttt{hyperref}, \texttt{listings}, \texttt{geometry},
\texttt{siunitx}, \texttt{enumitem}.
\end{itemize}

Sin LaTeX se pueden igual ejecutar tests y figuras; fallarán solo
\texttt{make paper}, \texttt{make docs} y \texttt{make guides}.

En Arch Linux, un conjunto suficiente es el metapaquete de TeX Live más
los estilos anteriores. En Debian/Ubuntu suele bastar
\texttt{texlive-latex-extra} y \texttt{texlive-lang-spanish}.

\subsection{Opcional}

\begin{itemize}
\item Instalar el paquete en modo editable:
\texttt{pip install -e ".[dev]"} (entonces no hace falta el
\texttt{sys.path} de los scripts, aunque ellos lo siguen insertando).
\item Git: \texttt{provenance.py} intenta guardar el hash de \texttt{HEAD}
en el JSON de cada figura; si no hay repo, el campo queda vacío.
\end{itemize}

\section{Cómo ejecutar el proyecto}

Sustituya \texttt{analisis\_numerico} por la ruta real del clon.

\subsection{1. Entorno}

\begin{lstlisting}
cd analisis_numerico
python3 -m venv .venv
source .venv/bin/activate
pip install -r requirements.txt
\end{lstlisting}

Compruebe \texttt{python3 --version} \(\ge\) 3.10. Si \texttt{python3} no
existe, use el intérprete concreto (por ejemplo \texttt{python3.12}).

\subsection{2. Tests (recomendado primero)}

\begin{lstlisting}
.venv/bin/pytest
# equivalente:
make tests
\end{lstlisting}

Debe terminar con todos los casos en verde salvo tres \texttt{skipped}.

\subsection{3. Figuras}

\begin{lstlisting}
.venv/bin/python scripts/tmm/run_all.py         # TMM, segundos
.venv/bin/python scripts/tmm/figure_04.py       # una sola
PWE_WORKERS=8 .venv/bin/python scripts/pwe/run_all.py   # PWE, horas
.venv/bin/python scripts/comparison/run_all.py  # requiere los dos
make results    # los tres
\end{lstlisting}

Salida esperada: bloques ``Figura N escrita\ldots'' y archivos en
\texttt{results/<método>/figure\_0N/output/}.

\subsection{4. Documentos PDF}

\begin{lstlisting}
make paper      # paper_reimplementation.pdf
make docs       # software_documentation.pdf
make guides     # guia_conceptual.pdf, solucion_numerica.pdf,
                # arquitectura.pdf
make pwe-docs   # metodo_ondas_planas.pdf, comparacion_tmm_pwe.pdf
make documents  # todos
make verify     # comprobacion final
\end{lstlisting}

\texttt{make guides} exige que las figuras TMM existan, y \texttt{make pwe-docs}
que existan \texttt{results/pwe/} y \texttt{results/comparison/}.

\subsection{5. Todo de una vez}

\begin{lstlisting}
source .venv/bin/activate
make all
\end{lstlisting}

Orden interno: tests \(\to\) resultados (TMM, PWE, comparación) \(\to\)
documentos \(\to\) verificación. Si los tests fallan, \texttt{make} se detiene
y no genera figuras.

\subsection{Cambiar un parámetro y volver a correr}

Edite el YAML (por ejemplo \texttt{omega\_m\_over\_2pi\_ghz} en
\texttt{configs/figure\_03.yaml}) y ejecute solo el script de esa figura.
No recompile el paper si no necesita el PDF actualizado.

Para TM: ponga \texttt{polarization: TM} en el YAML. El original no grafica
TM; el código sí la implementa (\(\chi=\varepsilon\)).

\subsection{Limpieza}

\begin{lstlisting}
make clean    # auxiliares de LaTeX de docs/ y presentacion/
\end{lstlisting}

No borra figuras ni el entorno virtual.

\section{Problemas frecuentes}

\begin{itemize}
\item \texttt{ModuleNotFoundError: fibonacci\_photonics}: está ejecutando un
script fuera de \texttt{scripts/} sin \texttt{src} en el path. Use los
scripts del repo o \texttt{PYTHONPATH=src}.
\item \texttt{ZeroDivisionError} en \(\omega=0\): el Drude diverge; el YAML
no debe incluir el origen en el barrido (la Fig.~1 empieza en
\(0.15\,\mathrm{GHz}\)).
\item Fig.~1 sin gaps a \(\theta=0\) y \(\omega_e=\omega_m\): no es un bug,
es adaptación de impedancia.
\item \texttt{pdflatex: command not found}: instale TeX Live; tests y
figuras no lo necesitan.
\item Figuras viejas en el paper: regenerar con
\texttt{scripts/tmm/run\_all.py} y luego \texttt{make paper}.
\item El PWE tarda mucho o satura la máquina: limite los procesos con
\texttt{PWE\_WORKERS=n}.
\item El \texttt{Makefile} apunta a \texttt{.venv/bin/python}. Si el venv
está en otra ruta, invoque los comandos con el intérprete explícito.
\end{itemize}

\section{Dónde está cada documento}

\begin{center}
\begin{tabular}{ll}
\toprule
PDF en la raíz & Contenido \\
\midrule
\texttt{paper\_reimplementation.pdf} & artículo de la reproducción \\
\texttt{software\_documentation.pdf} & manual corto de la API \\
\texttt{guia\_conceptual.pdf} & física desde cero \\
\texttt{solucion\_numerica.pdf} & TMM y código de las figuras \\
\texttt{arquitectura.pdf} & este documento \\
\texttt{metodo\_ondas\_planas.pdf} & formulación y resultados del PWE \\
\texttt{comparacion\_tmm\_pwe.pdf} & anexo de comparación TMM vs PWE \\
\bottomrule
\end{tabular}
\end{center}

Fuentes LaTeX en \texttt{docs/}.
Auditoría previa al código: \texttt{docs/fase\_A\_auditoria\_matematica.md}.
Informe de validación: \texttt{docs/INFORME\_VALIDACION.md}.
Paper original: PDF en la raíz del repo (no versionado).
Capítulo del libro y cuaderno del PWE: \texttt{references/pwe/}.

\section{Resumen operativo}

\begin{lstlisting}
python3 -m venv .venv && source .venv/bin/activate
pip install -r requirements.txt
pytest
python scripts/tmm/run_all.py
make paper docs guides
\end{lstlisting}

Eso deja tests en verde, las seis figuras y todos los PDF en la raíz del
proyecto.

\end{document}
